\documentclass[10pt,a4paper]{amsart}
\usepackage[margin=19mm]{geometry}
\usepackage{amsmath,amssymb,mathtools}
\usepackage[scr=rsfs,cal=euler]{mathalfa}
\usepackage{xfrac}
\usepackage[unicode,hidelinks,bookmarks=false]{hyperref}
\newcommand{\ie}{i.e.\ }
\newcommand{\cf}{cf.\ }
\newcommand{\resp}{respectively\ }
\newcommand{\Subsection}{\subsection}
\providecommand{\nobreakdash}{\nobreak\hbox{-}}
\setlength{\emergencystretch}{2em}
\multlinegap0pt
\usepackage{amssymb}
\newtheorem{theorem}{Theorem}
\newcommand{\fr}{\frac}
\title{Two-color partitions with evens in one color}
\author{George E. Andrews, Mohamed El Bachraoui}
\date{Source: arXiv:2512.23391v2 (CC BY 4.0)}
\begin{document}
\maketitle

\noindent\emph{Source and licence.} Two-color partitions with evens in one color. Version arXiv:2512.23391v2 (CC BY 4.0), distributed by arXiv under the Creative Commons Attribution 4.0 International licence, http://creativecommons.org/licenses/by/4.0/, as recorded in the arXiv licence metadata for this exact version and shown on the abstract page https://arxiv.org/abs/2512.23391v2. Full licence notice: \texttt{licenses/arxiv-2512.23391-CC-BY-4.0.txt}.\par\smallskip

\noindent\textit{Editorial scope.} Counted unit: the article's theorem F01-id (source0.tex lines 216--222). The article's complete original proof of it is delivered with it (source0.tex lines 224--245, one \texttt{proof} environment of 22 lines). No mathematics is written, repaired or re-derived: the statement and the proof are the article's own lines, in the article's own notation, and no retained line is substituted or re-flowed.  Retained uncounted as the source-local context the proof needs: lines 192--194; lines 196--200.  Nothing was omitted from any retained span, so the package carries no omission record.  No retained line contains a citation, so the wrapper carries no bibliography and no bibliography entry.  No retained line contains a figure environment, an image-inclusion command or drawing code, so no image is delivered and the package declares no asset dependency.  Editorial formatting only, in addition: an AMS \texttt{amsart} preamble that restates the article's own declarations and macros used by the retained lines, namely the environments \texttt{theorem} on separate counters, together with the standard packages those lines need.  The retained environments print in this document as Theorem 1 in order of appearance, whereas the article prints the same items with the numbering of its own full document.  The complete native source of the article as distributed in the arXiv e-print for this exact version is bundled unchanged as \texttt{sources/191803/source0.tex}.
\begin{equation}\label{Euler}
(-q;q)_\infty =\fr{1}{(q;q^2)_\infty},
\end{equation}

\begin{align}
\sum_{n\geq 0} F(n) q^n &=\fr{1}{(q;q^2)_\infty^2 (q^2;q^2)_\infty} =\fr{(-q;q)_\infty}{(q;q^2)_\infty (q^2;q^2)_\infty}
= \fr{(-q;q)_\infty}{(q;q)_\infty} \nonumber \\
&=\sum_{n\geq 0}\overline{p}(n) q^n, \label{F-id}
\end{align}

\begin{theorem}\label{thm F01-id}
For any nonnegative integer $n$ there holds
\begin{align}
F_0(n) &= \fr{\overline{p}(n)+\overline{p}\big(\fr{n}{2}\big)}{2}, \label{F0-id} \\
F_1(n) &= \fr{\overline{p}(n)-\overline{p}\big(\fr{n}{2}\big)}{2}, \label{F1-id}.
\end{align}
\end{theorem}

\begin{proof}
It is easily seen that
\[
\sum_{n\geq 0} \big( F_0(n)-F_1(n) \big) q^n =\fr{1}{(-q;q^2)_\infty (q;q^2)_\infty (q^2;q^2)_\infty}.
\]
Then by~\eqref{Euler} and~\eqref{F-id}, we obtain
\begin{align}
\sum_{n\geq 0} \big( F_0(n)-F_1(n) \big) q^n &=\fr{1}{(-q;q^2)_\infty (q;q^2)_\infty (q^2;q^2)_\infty} \nonumber \\
&=\fr{(-q;q)_\infty}{(-q;q^2)_\infty (q^2;q^2)_\infty}  \nonumber\\
&=\fr{(-q^2;q^2)_\infty}{(q^2;q^2)_\infty} \nonumber \\
&=\sum_{n\geq 0} \overline{p}(n) q^{2n} . \label{G0-G1}
\end{align}
Furthermore, we clearly have
\[
\sum_{n\geq 0} \big( F_0(n)+F_1(n) \big) q^n = \sum_{n\geq 0} F(n) q^n,
\]
and thus by~\eqref{F-id},
\begin{equation}\label{gen G0pG1}
\sum_{n\geq 0} \big( F_0(n)+F_1(n) \big) q^n = \sum_{n\geq 0}\overline{p}(n) q^n.
\end{equation}
Now combine~\eqref{G0-G1} with~\eqref{gen G0pG1} to achieve the desired formulas.
\end{proof}

\end{document}
